%% ma: L12-B08-0168
%% lop: 12
%% chuong: 2
%% bai: 8
%% dang: 12.8.53
%% loai: TN4
%% muc: VD
%% dap_an: "B"
%% nguon: "CĐ2.4.pdf giáo viên gửi 27/10/2026 (M07, Lê Minh Khôi – Nguyễn Anh Khôi, Chương II Chủ đề 4 Các bài toán về tâm tỉ cự), Câu 18"
%% kien_thuc: "Bài 8 (khoảng cách hai điểm, biểu thức tọa độ), Bài 6 (bình phương vô hướng, tâm tỉ cự)"
%% trang_thai: cho-duyet
%% ly_do: "Dạng mới 12.8.53 đề xuất, chờ giáo viên duyệt"
\begin{ex}
Trong không gian $Oxyz$, cho ba điểm $A(0;0;-1)$, $B(-1;1;0)$, $C(1;0;1)$. Tọa độ của điểm $M$ thỏa mãn $3MA^2+2MB^2-MC^2$ đạt giá trị nhỏ nhất là
\choice{$M\left(\dfrac34;-\dfrac12;1\right)$}{\True $M\left(-\dfrac34;\dfrac12;-1\right)$}{$M\left(-\dfrac34;-\dfrac12;1\right)$}{$M\left(\dfrac34;\dfrac12;-1\right)$}
\loigiai{Gọi $I$ thỏa $3\overrightarrow{IA}+2\overrightarrow{IB}-\overrightarrow{IC}=\vec 0\Rightarrow I=\dfrac{3A+2B-C}4=\left(-\dfrac34;\dfrac12;-1\right)$.

$3MA^2+2MB^2-MC^2=4MI^2+3IA^2+2IB^2-IC^2$, nhỏ nhất khi $M\equiv I$.}
\end{ex}
